% beam_comparison_text.tex -- GENERATED, do not edit by hand.
%
% Draft text for the antenna comparison. Regenerate by running
% horizon_position/notebooks/beam_comparison.ipynb in the mock_analysis repo;
% it is generated from beam_comparison_text.tex.in by substituting the
% placeholders. Every number is computed at generation time from the arrays
% the figure is drawn from, for the reason set out in signal_loss_text.tex.in:
% the one block that ever carried hand-entered numbers is the one whose
% numbers went stale.
%
% WHAT THIS IS FOR. The manuscript already claims, in
% \subsection{subsec:eig_antenna}, that the bowtie "design is chosen to
% minimize the overlap of the spectral shapes of the beam pattern and
% reflection coefficient with models of the 21 cm signal". Until this figure
% that claim was asserted and never measured, and Fig.~\ref{fig:singular_values}
% quantified the overlap for exactly one antenna, so there was nothing to
% compare it against. That is what made the signal-loss numbers read as
% "mediocre" rather than "good": mediocre relative to what?
%
% FRAMING, and this matters more than the numbers. The Vivaldi is NOT a
% rejected candidate. It is a HERA Phase II feed (\citealt{2021ITAP...69.8143F},
% \citealt{2024PASP..136d5002B}), it is what EIGSEP uses for the ground
% antennas (section~\ref{sec:vivaldis}), and it is the antenna the first
% suspension flew in October 2024 (the paragraph beginning "The suspension with
% the highline and pulley system was first used..."). So the honest framing is
% a documented evolution -- the first suspension flew this feed, the bowtie was
% built afterwards for this job, here is what that bought -- and NOT a
% retrospective justification of a choice between two candidates.
%
% Two things the text must keep doing or this backfires:
%   1. Say the feed is used outside the configuration it was optimised for.
%      Section~\ref{sec:vivaldis} already states EIGSEP uses it without the
%      14-m dish it was designed to illuminate; this text leans on that. Do not
%      let any sentence read as a criticism of the HERA design.
%   2. Head off "are the ground antennas compromised?". A referee will notice
%      the ground antennas are Vivaldis. One clause on their differing role
%      closes it. Do not drop that clause to save a line.
%
% THE VIVALDI COUPLES BETTER TO THE SKY AND STILL LOSES. Keep this. It is the
% single most persuasive sentence available, because the alternative antenna is
% handed an advantage and still comes out behind: a one-sided comparison reads
% as special pleading, this does not. ETAVIV > ETABOW is asserted in the
% notebook so the sentence cannot outlive the fact.
%
% RETAINED FRACTION, NOT ABSOLUTE MILLIKELVIN. The comparison is quoted as a
% fraction of the unfiltered signal at MATCHED foreground suppression. Both
% halves are load-bearing. Fraction, because absolute retained amplitude mixes
% in aperture efficiency and the two antennas are then nearly tied -- that is a
% statement about collecting area, not about spectral overlap, which is what
% the section is about. Matched suppression, because at a fixed MODE COUNT the
% Vivaldi appears to retain MORE, for the trivial reason that it has suppressed
% the foregrounds less; a panel plotting retention against N inverts the
% conclusion and a referee would find it immediately.
%
% NO MODE COUNT IS AN OPERATING POINT, here either. NABOW/NAVIV/NAISO are where
% each antenna's foreground residual drops below its own median retained
% signal, and the figure marks none of them. Same rule as
% signal_loss_text.tex.in: no threshold is drawn, no depth is adopted.
% ===================================================================


% ===================================================================
% BLOCK 1 -- caption for beam_comparison.pdf. Two-column `figure*'; the row is
% 7in wide.
%
% The caption must say that the grey bands DIFFER between panels. Each antenna
% has its own open-sky fraction and its own eigenbasis, so the same models
% project differently; a reader who assumes it is one band drawn three times
% makes exactly the error a single-panel version of this figure would have
% forced. It must also say the isotropic residual leaves the frame, or the
% truncated curve reads as missing data.
%
% What it must NOT do is explain WHY the bands differ beyond naming the cause.
% The three-panel layout carries the point; a sentence arguing it invites the
% argument.
% ===================================================================

\caption{Residual RMS of the simulated antenna temperature (black) and of the
1769 global 21-cm models filtered with the same modes (5--95 per cent of
ensemble in grey band, median dashed), as a function of the number of
eigenmodes filtered ($N$), for three antennas observing the same sky behind the
same horizon. \textit{Left:} an isotropic beam, which has no chromaticity and
so isolates what the sky alone contributes. \textit{Centre:} the EIGSEP bowtie,
reproducing Fig.~\ref{fig:singular_values}. \textit{Right:} a Vivaldi feed of
the type used for the ground antennas, without the dish it was designed to
illuminate. Each antenna is filtered in its own eigenbasis and its models are
scaled by its own beam-weighted open-sky fraction, so the grey bands differ
between panels. The isotropic residual falls below the axis near $N=6$,
where it reaches machine precision.}


% ===================================================================
% BLOCK 2 -- section "Minimising Covariance with the 21-cm Signal", after the
% block-1 passage of signal_loss_text.tex (i.e. after the paragraph ending
% "...as described in section~\ref{subsec:fwd_modelling}").
%
% Three paragraphs: what was compared, what it costs, what it means for the
% design. Results first; the caveats are one sentence each and they come last.
%
% Paragraph 3 is the destination -- the design claim in
% section~\ref{subsec:eig_antenna} is the thing this is here to substantiate.
% If the block has to be cut, cut towards that sentence.
% ===================================================================

The amount of signal that survives this filter is a property of the antenna as
much as of the sky. To separate the two, we repeated the simulation and the
filter for three beams observing the same sky behind the same horizon: the
EIGSEP bowtie, a Vivaldi feed of the type used for the ground antennas
(section~\ref{sec:vivaldis}) and flown on the first suspension test, and an
isotropic beam with no chromaticity at all, which bounds what any antenna could
achieve at this site. Each is filtered in its own eigenbasis, since the modes
are a property of the beam-weighted foregrounds it produces.

The antenna sets both how deep the filter must go and how much signal is left
when it gets there. Filtering the beam-weighted foregrounds to 1\,mK
takes 6 modes for the isotropic beam, 10 for the bowtie and
18 for the Vivaldi feed; of the unfiltered 21-cm RMS, 13,
8 and 4 per cent respectively remain at that point. Compared
at the same level of foreground suppression, the bowtie therefore retains
2.2 times the signal the Vivaldi feed does, and recovers 58 per
cent of what an achromatic antenna would. This is not an aperture effect and
runs against one: the Vivaldi feed is the more directive of the two and couples
proportionally more of its response to the sky than to the ground
(0.70 against 0.45), so it begins with more signal and still ends
with less.

This is the design criterion of section~\ref{subsec:eig_antenna} measured
rather than asserted, and it is why the suspended antenna is not simply a
Vivaldi feed of the kind already deployed on the ground. The comparison is
specific to this application: the feed is used here in isolation, without the
dish it was designed to illuminate, and the ground antennas it is used for are
not required to deliver a low-chromaticity global-signal measurement. As
throughout this section, no mode count above is a filter depth we propose to
apply to data; the figure marks none of them, and the quantity being compared
is spectral overlap, not sensitivity.
